% =================================================
% MAIN DOCUMENT (CLASSIC ARTICLE)
% =================================================
%
% Classic LaTeX article. Layout and formatting live in
% classicstyle.sty. Font, math packages, and macros
% are chosen here.
%
% =================================================


\documentclass[12pt]{article}


% =================================================
% LOAD STYLE
%
% Tools available via classicstyle.sty:
%   - figures:   graphicx, float ([H]-placement), wrapfig, subfig
%   - tables:    booktabs (\toprule, \midrule, \bottomrule), caption
%   - lists:     enumitem (custom labels and spacing)
%   - boxes:     tcolorbox (colored frames and notes)
%   - colors:    xcolor with svgnames
%   - citations: natbib (numerical, sorted and compressed)
%
% See README.md for syntax examples.
% =================================================
\usepackage{classicstyle}


% =================================================
% LANGUAGE
% =================================================
\usepackage[english]{babel}


% =================================================
% FONT
%
% Choose one text font. Comment out the others.
% =================================================

\usepackage{mlmodern}              % Modern Latin Modern
% \usepackage{lmodern}             % Latin Modern (classic)
% \usepackage{stix2}               % STIX2 (math and text)

% Sans-serif font (optional):
\usepackage{helvet}                % Helvetica as sans font

% For sans-serif headings, uncomment:
% \usepackage{titlesec}
% \titleformat*{\section}{\sffamily\Large\bfseries}
% \titleformat*{\subsection}{\sffamily\large\bfseries}


% =================================================
% MATHEMATICS
% =================================================
\usepackage{amsmath, amsfonts, amsthm}
\usepackage{sansmath}              % math in sans
\usepackage{siunitx}               % SI units
\usepackage{esvect}                % nice vectors


% =================================================
% CUSTOM PACKAGES
%
% Add project-specific packages here.
% =================================================

% \usepackage{tikz}                % drawings / diagrams
% \usepackage{physics}             % \dd, \ket, \bra, ...


% =================================================
% CUSTOM COMMANDS
%
% Personal mathematical shortcuts.
% =================================================

% --- Greek letters / constants ---
\newcommand{\p}{\varphi}
\newcommand{\ep}{\varepsilon}
\newcommand{\e}{\mathrm{e}}

% --- Unit vectors ---
\newcommand{\er}{\vec{e}_r}
\newcommand{\ephi}{\vec{e}_\varphi}

% --- Derivatives ---
\newcommand{\pdd}[2]{\frac{\partial #1}{\partial #2}}
\newcommand{\pddd}[1]{\frac{\partial^2}{\partial #1^2}}
\newcommand{\dd}[2]{\frac{\mathrm{d} #1}{\mathrm{d} #2}}
\newcommand{\diff}[1]{\mathrm{d}{#1}~}
\newcommand{\diffn}[2]{\mathrm{d}^{#1} {#2}~}
\newcommand{\diffnn}[2]{\mathrm{d}^{#2} {#1}~}

% --- Partial derivatives ---
\newcommand{\ddx}{\frac{\partial }{\partial x}}
\newcommand{\ddy}{\frac{\partial }{\partial y}}
\newcommand{\ddr}{\frac{\partial }{\partial r}}
\newcommand{\dddx}{\frac{\partial^2}{\partial x^2}}
\newcommand{\dddy}{\frac{\partial^2}{\partial y^2}}
\newcommand{\ddphi}{\frac{\partial }{\partial \varphi}}
\newcommand{\dddr}{\frac{\partial^2 }{\partial r^2}}
\newcommand{\dddphi}{\frac{\partial^2 }{\partial \varphi^2}}

% --- Integration ---
\newcommand{\pint}[2]{\int\!{#1}\,d{#2}}
\newcommand{\pintlim}[4]{\int_{#1}^{#2}\!{#3}\,d{#4}}

% --- Algebra ---
\newcommand{\nsqrt}[2]{\sqrt[\uproot{2}\leftroot{5} {#1}] {#2}}


% =================================================
% METADATA
% =================================================
\title{\huge Article Title}
\author{\small First Last}
\date{Spring 2026}


% =================================================
% DOCUMENT
% =================================================
\begin{document}

{\fontfamily{lmr}\selectfont\maketitle}     % title in Latin Modern
\vspace{1cm}

\begin{quote}
  \emph{``A classical quote goes here, setting the tone for the
  article that follows.''}
  \begin{flushright}
    --- Author
  \end{flushright}
\end{quote}

\begin{abstract}
  \noindent The abstract gives the reader a quick overview of what
  the article is about, what the central argument is, and why it is
  relevant. Two to four sentences is usually sufficient.
\end{abstract}

\newpage


\section{Introduction}

The introduction states the problem and places it in context.
Paragraphs are marked with indentation rather than line breaks,
so the text forms a continuous flow. This is the classical
typographic convention for printed articles and books.

A new paragraph continues with the same rhythm. References to prior
work are made with a citation~\cite{example2024}, and several sources
can be combined within the same parentheses~\cite{example2024,examplebook2020}.


\section{Main Results}

Mathematics is set either inline, as $E = mc^2$, or as a numbered
equation:
\begin{equation}
  \hat{H}\psi = E\psi.
  \label{eq:example}
\end{equation}
Reference equations with \eqref{eq:example} where it fits naturally
in the text.

A new paragraph after an equation begins with indentation, just as
elsewhere. The spacing between equation and body text is set to
0.5 cm in \texttt{classicstyle.sty}, so the mathematics breathes
without detaching from the argument.

\begin{figure}[H]
  \centering
  \includegraphics[width=0.7\linewidth]{Energy_levels}
  \caption{An example figure. The \texttt{[H]} placement locks the
           figure exactly where it sits in the source, while \texttt{[t]}
           lets it float to the top of the page.}
  \label{fig:example}
\end{figure}

Refer to figures in the text with figure~\ref{fig:example}. To
summarize key values, tables can be useful:

\begin{table}[H]
  \centering
  \caption{Example of a table set with \texttt{booktabs}.}
  \label{tab:example}
  \begin{tabular}{lcr}
    \toprule
    Quantity            & Symbol  & Value                         \\
    \midrule
    Speed of light      & $c$     & $2.998 \times 10^8$ m/s       \\
    Planck's constant   & $\hbar$ & $1.055 \times 10^{-34}$ Js    \\
    Elementary charge   & $e$     & $1.602 \times 10^{-19}$ C     \\
    \bottomrule
  \end{tabular}
\end{table}

Vertical lines are deliberately omitted, and horizontal lines are
used sparingly --- only top, middle, and bottom. This gives the
table a calm, professional look.


\section{Discussion}

The discussion is where you interpret the results and place them
in the context of existing literature. Here you can be more
speculative than in the results section, as long as it is clear
which arguments are your own and which are established findings.

Bullet lists are useful when several equally weighted points
should be enumerated:
\begin{itemize}
  \item The first point is stated here.
  \item A new point follows on the next line.
  \item A third point closes the list.
\end{itemize}

For numbered sequences or step-by-step reasoning, use numbered
lists:
\begin{enumerate}
  \item The first step in the reasoning.
  \item The second step, building on the first.
  \item The third step leads to the conclusion.
\end{enumerate}

\begin{tcolorbox}[colback=Beige, colframe=DarkSlateGray, title=Note]
  Colored boxes are useful for highlighting definitions,
  observations, or warnings. The \texttt{tcolorbox} package
  offers many design options --- see its documentation for
  the full set.
\end{tcolorbox}

\subsection{A Subsection}

Subsections get their own heading but keep a calm typographic
weight, so the main sections still dominate the hierarchy.


\section{Conclusion}

The conclusion summarizes briefly, points toward open questions,
and suggests further work where relevant. A final paragraph rounds
off the article without introducing new arguments.


% =================================================
% REFERENCES
% Compile order:
%   pdflatex main
%   bibtex   main
%   pdflatex main
%   pdflatex main
% =================================================
\bibliographystyle{unsrtnat}
\bibliography{references}


\end{document}
